\subsection{Completely integrally closed domains}

DEFINITION 5. An integral domain A is called completely integrally closed if the following condition is satisfied: every element x of the field of fractions K of A such that all the powers \(x^{n}\) ( \(n \geqslant 0\) ) are contained in a finitely generated sub-A-module of K, belongs to A.

Note that the hypothesis that the \(x^n\) are contained in a finitely generated sub-A-module of K can also be expressed by saying that there exists a non-zero element \(d \in \mathbf{A}\) such that \(dx^n \in \mathbf{A}\) for all \(n \geqslant 0\) ; for the latter condition means that \(x^n \in \mathrm{Ad}^{-1}\) ; and conversely, if \((b_i)_{1 \leqslant i \leqslant m}\) is a finite sequence of elements of K, there exists \(d \in \mathbf{A}\) such that \(db_i \in \mathbf{A}\) for \(1 \leqslant i \leqslant m\) , whence \(d\mathbf{M} \subset \mathbf{A}\) for the sub-A-module M of K generated by the \(b_i\) .

Clearly a completely integrally closed domain is integrally closed; conversely, the Corollary to Proposition 1 of no. 1 shows that an integrally closed Noetherian domain is completely integrally closed. * On the other hand, the ring of a valuation of height ≥ 2 (Chapter VI, § 4, no. 4) is integrally closed but not completely integrally closed. * If (A\_t) is a family of completely integrally closed domains with the same field of fractions K, A = ∩\_t A\_t is completely integrally closed. For if x ∈ K is such that for some non-zero d in A, \(dx^n\) belongs to A for all n \textgreater{} 0, the hypothesis implies that x ∈ A\_t for all t and hence x ∈ A.

PROPOSITION 14. Let A be a completely integrally closed domain. Then every polynomial ring \(\mathbf{A}[\mathbf{X}_1,\ldots ,\mathbf{X}_n]\) (resp. every ring of formal power series \(\mathbf{A}[[\mathbf{X}_1,\dots ,\mathbf{X}_n]]\) ) is completely integrally closed.

By induction on \(n\) , it is sufficient to prove that A{[}X{]} (resp. A{[}{[}X{]}{]}) is completely integrally closed. Then let P be an element of the field of fractions of A{[}X{]} (resp. A{[}{[}X{]}{]}) and suppose that there exists a non-zero element Q ∈ A{[}X{]} (resp. Q ∈ A{[}{[}X{]}{]}) such that QP \(^{m}\) ∈ A{[}X{]} (resp. QP \(^{m}\) ∈ A{[}{[}X{]}{]}) for every integer m ≥ 0. If K is the field of fractions of A, A{[}X{]} (resp. A{[}{[}X{]}{]}) is a subring of K{[}X{]} (resp. K{[}{[}X{]}{]}) and K{[}X{]} (resp. K{[}{[}X{]}{]}) is a principal ideal domain (Algebra, Chapter VII, § 1, no. 1) and hence integrally closed (no. 4, Proposition 10) and Noetherian (Algebra, Chapter VIII, § 2, no. 3) and therefore completely integrally closed; then we have already seen that \(\mathbf{P} \in \mathbf{K}[\mathbf{X}]\) (resp. \(\mathbf{P} \in \mathbf{K}[[\mathbf{X}]]\) ). Let \(\mathbf{P} = \sum_{k=0}^{\infty} a_k \mathbf{X}^k\) ( \(a_k \in \mathbf{K}\) ) and \(\mathbf{Q} = \sum_{k=0}^{\infty} b_k \mathbf{X}^k\) ( \(b_k \in \mathbf{A}\) ) and we argue by reductio ad absurdum by supposing that the \(a_k\) do not all belong to \(\mathbf{A}\) ; then there is a least index \(i\) such that \(a_i \notin \mathbf{A}\) ; if we write

\(P_{1}=\sum_{k=0}a_{k}X^{k}\in A[X]\) , it follows immediately from the hypothesis that also \(Q(P-P_{1})^{m}\in A[X]\) (resp. \(Q(P-P_{1})^{m}\in A[[X]]\) ) for all \(m\geqslant0\) . Let j be the least integer such that \(b_{j}\neq0\) ; clearly in \(Q(P-P_{1})^{m}\) the term of least degree with a coefficient \(\neq0\) is \(b_{j}a_{i}^{m}X^{j+m^{i}}\) and hence \(b_{j}a_{i}^{m}\in A\) for all \(m\geqslant0\) ; but as A is completely integrally closed this implies \(a_{i}\in A\) , contrary to the hypothesis.

PROPOSITION 15. Let A be a filtered ring whose filtration is exhaustive and such that every principal ideal of A is closed under the topology defined by the filtration. If the associated graded ring gr(A) (Chapter III, §2, no. 3) is a completely integrally closed domain, then A is a completely integrally closed domain.

Let \((\mathbf{A}_n)_{n\in \mathbf{Z}}\) be the filtration defined on \(\mathbf{A}\) ; as \(\bigcap_{n\in \mathbf{Z}}\mathbf{A}_n\) is the closure of the ideal (0) (Chapter III, § 2, no. 5), the hypothesis implies first that the filtration \((\mathbf{A}_n)\) is separated and, as \(\operatorname{gr}(\mathbf{A})\) is an integral domain, so then is \(\mathbf{A}\) (Chapter III, § 2, no. 3, Corollary to Proposition 1). Let \(x = b / a\) be an element of the field of fractions \(\mathbf{K}\) of \(\mathbf{A}\) ( \(a\in \mathbf{A}, b\in \mathbf{A}\) ) for which there exists an element \(d\neq 0\) of \(\mathbf{A}\) such that \(dx^n\in \mathbf{A}\) for all \(n\geqslant 0\) . We must prove that \(b\in \mathbf{A}a\) and, as by hypothesis the ideal \(\mathbf{A}a\) is closed, it is sufficient to show that, for all \(n\in \mathbf{Z}, b\in \mathbf{A}a + \mathbf{A}_n\) . As the filtration of \(\mathbf{A}\) is exhaustive, there exists an integer \(q\in \mathbf{Z}\) such that \(b\in \mathbf{A}a + \mathbf{A}_q\) . It will therefore suffice to prove that the relation \(b\in \mathbf{A}a + \mathbf{A}_m\) implies \(b\in \mathbf{A}a + \mathbf{A}_{m + 1}\) .

Suppose then that \(b = ay + z\) where \(y \in \mathbf{A}\) , \(z \in \mathbf{A}_m\) . Then by hypothesis \(dx^n \in \mathbf{A}\) for all \(n \geqslant 0\) , whence we obtain immediately \(d(x - y)^n \in \mathbf{A}\) for all \(n \geqslant 0\) ; in other words, \(dz^n = a^n t_n\) where \(t_n \in \mathbf{A}\) for all \(n \geqslant 0\) . We can obviously limit our attention to the case where \(z \neq 0\) . Let \(v\) denote the order function on \(\mathbf{A}\) (Chapter III, §1, no. 2) and let us write \(v(d) = n_1\) , \(v(z) = n_2 \geqslant m\) , \(v(a) = n_3\) ; let \(d', z', a'\) be the respective images of \(d, z, a\) in \(\mathbf{A}_{n_1}/\mathbf{A}_{n_1+1}\) , \(\mathbf{A}_{n_2}/\mathbf{A}_{n_2+1}\) , \(\mathbf{A}_{n_3}/\mathbf{A}_{n_3+1}\) . For all \(n \geqslant 0\) , \(v(dz^n) = n_1 + nn_2\) (Chapter III, §2, no. 3, Proposition 1) and hence the canonical image in \(\mathrm{gr}(\mathbf{A})\) of \(dz^n\) is \(d'z'^n\) ; similarly it is seen that the canonical image in \(\mathrm{gr}(\mathbf{A})\) of \(a^n t_n\) is of the form \(a'^n t_n'\) where \(t_n' \in \mathrm{gr}(\mathbf{A})\) and, as \(a' \neq 0\) we deduce that, for all \(n \geqslant 0\) , \(d'(z'/a')^n \in \mathrm{gr}(\mathbf{A})\) . The hypothesis that \(\mathrm{gr}(\mathbf{A})\) is completely integrally closed therefore implies the existence of an \(s' \in \mathrm{gr}(\mathbf{A})\) such that \(z' = a's'\) ; decomposing \(s'\) into a sum of homogeneous elements, it is further seen (since \(z'\) and \(a'\) are homogeneous) that \(s'\) may be assumed to be homogeneous, that is the image of an element \(s \in \mathbf{A}\) ; then \(v(as) = v(z) = n_2\) and \(z \equiv as (\mathrm{mod. A_{n_2+1}})\) ; as \(n_2 \geqslant m\) , a fortiori \(z \equiv as (\mathrm{mod. A_{m+1}})\) , hence \(b \equiv a(y + s) (\mathrm{mod. A_{m+1}})\) .
% semantic-fragments: legacy-input-seam
\relax{}
